\section{Economic Policy}

\begin{frame}{Shell Model}
    \mycomment{
    There are many levels of government, from the central government over federal states down to the communal level; and on each level, there are multiple ministries, government owned companies, and institutions.

    These sometimes complement, and sometimes compete. It is worthwhile and not an easy task to take a more differentiated view. In particular, government statistics need to be carefully read, for example on budget and debt level of government.
    }
    \showmytikz{../shared/tikz/tikz-shellmodel}
\end{frame}

\begin{frame}{Objectives of Macroeconomic Policy}
    \showmytikz[1]{../shared/tikz/tikz-economicsmagicsquare}
    \mycomment{
    The so-called \enquote{magic square} (\textit{Magisches Viereck}) is a popular design in the German discourse. This is commonly understood to be \enquote{gesamtwirtschaftliches Gleichgewicht} as required per GG Art. 109 / Par. 1 StabG.

    In the English literature, this is often known as \enquote{Kaldor's Diamond}, though Kaldor's publication came later than the German concept.

      }
\end{frame}

\begin{frame}{Objectives of Macroeconomic Policy}
    \showmytikz[2]{../shared/tikz/tikz-economicsmagicsquare}
    \mycomment{
    There has been proposals to prefer other objectives (Sebastian \noun{Dullien} et al).

    In both diagrams, it remains questionable if, or to which extent, these goals compete with one another, or if they can indeed only be achieved in unison.
    }
\end{frame}

\begin{frame}{Objective of Macroconomic Policy - EU}
    \showmytikz{../shared/tikz/tikz-119TFEU}
    \mycomment{There is a similar concept in Art. 119 TFEU on the EU level. 
    
    }
\end{frame}

\begin{frame}{Objective of Macroconomic Policy - EU}
    \showmytikz{../shared/tikz/tikz-3TEU}
    \mycomment{
    Art. 3 TEU defines the basis and the aims of the internal market of the Union.
    }
\end{frame}

\begin{frame}{Other Objectives}
    \mytext{
    Correction of market failures
    \begin{itemize}
        \item Unemployment insurance
        \item Health insurance
        \item Monopolies
        \item Poverty traps
        \item Infrastructure
        \item Education
    \end{itemize}
    }
\end{frame}

\begin{frame}{Doughnut Economics}
    \showmytikz{../shared/tikz/tikz-doughnut}
    \mycomment{
    A recent concept by Kate \noun{Raworth}. The core idea is that economic activity should have, or has, corner conditions or boundaries. 
    }
    
\end{frame}

\begin{frame}{Sustainable Development Goals, UN}
    \showmytikz{../shared/tikz/tikz-sdggrid}
    \mycomment{
    The 17 SDGs have been criticised. Why these 17? And how are they quantifiable? Are these actually goals, or areas of activity?
    }
\end{frame}

\begin{frame}{Market Conformity of Economic Policy}
    \mytext{
    Economic Policy should take the insights from Microeconomics into account:
    \begin{itemize}
        \item prices depend on volumes, and volumes on prices
        \item there is always a dead weight loss
        \item there is always an excess burden
        \item there is always crowding out
        \item there are income and substitution effects
        \item ...
    \end{itemize}
    }
\end{frame}

\begin{frame}{Market Conformity of Economic Policy}
    \mytext{
    The macroeconomic perspective adds to that:
    \begin{itemize}
        \item an invention in one market has an influence on (all) other markets
        \item again, there are income and substitution effects
        \item thus, there are macroeconomic effects: every policy (of a certain size) affects interest rates, growth, employment, etc.
        \item micro-policy (into individual markets) is different to macro-policy 
        \begin{itemize}
            \item spending on one market: micro-policy
            \item government spending as such: macro-policy
        \end{itemize}
        \item micro- and macro-policy may conflict
    \end{itemize}

    }
\end{frame}


\begin{frame}{Order vs. Process Policy}
    \mycomment{
    Economic policy can take a multitude of forms, which can be classified in a number of different ways. No one scheme is exhaustive.

    One major distinction is that between order policy, that is the government setting the framework for liberal economic activity, and process policy, which influences what is happening on markets directly.
    }
    \showmytree[1]{../shared/tikz/forest-economic-policy}
\end{frame}

\begin{frame}{Fiscal Policy}
    \mycomment{
    Within process policy, any policy that touches the government budget is fiscal policy. This includes the income side of the government: tax policy; and the expenditure side. 
    
    The government has a plethora of tools at its disposal. These can be more broadly applied, or more specific and discreet. The government policy itself can be more or less rule based, depending on the constitutional and legal situation.
    }
    \showmytree[only 2]{../shared/tikz/forest-economic-policy}
\end{frame}

\begin{frame}{Further Means of Economic Policy}
    \mycomment{
    Apart from fiscal policy, the government has other means to influence economic activities. One important branch is monetary policy, which resides with the Central Bank. This we will cover extensively later.

    Structural policy overlaps with fiscal policy; apart from the immediate fiscal effects, policies have spacial and structural effects. One important tool is the establishment of special economic areas, where incentives to invest are set by tax rebates, lower standards, or subsidies.

    The government can also promote industries or investment in certain areas or sectors by providing information services. And trade agreements and foreign policy also impacts economic activities.
    }
    \showmytree[3]{../shared/tikz/forest-economic-policy}
\end{frame}

\begin{frame}{Why Economic Policy?}
    \mytext{
        We know from Microeconomics that markets are Pareto-efficient.
        \begin{itemize}
            \item Why \textit{should} a government pursue economic policy?
            \item[$\rightarrow$] normative
            \item Why \textit{would} a government pursue economic policy?
            \item[$\rightarrow$] descriptive
        \end{itemize}        
        Economic Policy is intervention!
    }
\end{frame}

\begin{frame}{Normative Arguments for Intervention}
    \mytext{
        \begin{itemize}
            \item Correct market failure
            \begin{itemize}
                \item under-use of resources
                \item inequality in chances
                \item over-use of resources
                \item externalities, tragedy of the commons
            \end{itemize}
            \item The government may be wiser than the markets
            \begin{itemize}
                \item long-term strategy
                \item more information than subjects
                \item morally superior 
            \end{itemize}
            \item government can raise revenue
        \end{itemize}
    }
\end{frame}

\begin{frame}{Normative Arguments against Intervention}
    \mytext{
    \begin{itemize}
        \item government not wiser than markets
        \item populism leads to irrational policy
        \item chauvinism, favouritism, nepotism
        \item tends to be conservative rather than innovative
        \item self-interest of government
        \item interest of voters: median voter theorem
        \item excess burden, crowding out, dead weight loss
        \item over-use, low prices, little investment
        \item rent-extraction, corruption
        \item debt financing burdens future generations
        \item unsustainable debt: frequent crises
        \item ...
    \end{itemize}
    }
\end{frame}